% !TeX program = xelatex
% UTF-8; standalone or input into a document using XePersian.
% Transcription of Notebook/altered mental status, handwritten pages 1--6.
% pic01.jpg = 20261013_065539198.jpg
% pic02.jpg = 20261013_065545726.jpg
% pic03.jpg = 20261013_065549914.jpg
% pic04.jpg = 20261013_065555081.jpg
% pic05.jpg = 20261013_065559654.jpg
% pic06.jpg = 20261013_065603179.jpg
% Medical statements and values are reproduced from the source, not corrected.
\ifdefined\settextfont
  \def\AMSStandalone{0}
\else
  \def\AMSStandalone{1}
  \documentclass[11pt,a4paper]{article}
  \usepackage[margin=2cm]{geometry}
  \usepackage{amsmath,amssymb,array,xcolor,enumitem}
  \usepackage{xepersian}
  \IfFontExistsTF{Amiri}{\settextfont{Amiri}}{\settextfont{FreeSerif}}
  \setlatintextfont{FreeSerif}
  \setlist[itemize]{itemsep=3pt,topsep=4pt,parsep=0pt}
\fi
\providecommand{\en}[1]{\lr{#1}}
\newcommand{\amsarr}{\ensuremath{\leftarrow}}
\newcommand{\amshl}[1]{\ifdefined\colorbox\colorbox{yellow!35}{\strut #1}\else\textbf{#1}\fi}
\newcommand{\amspage}[2]{%
  \clearpage\section*{\en{184 — Altered Mental Status} — صفحهٔ #1}
  {\footnotesize\noindent\en{Subject: 184 — AMS}\hfill\en{Date: / /}\par
  \noindent\en{#2}\par}\medskip}
\ifnum\AMSStandalone=1 \begin{document}\fi

\amspage{۱}{pic01.jpg}
\subsection*{\en{Consciousness / Arousal}}
\begin{latin}
\begin{itemize}
  \item \amshl{Consciousness} $\rightarrow$ \amshl{Cortical} function:
    awareness of self and environment (place, time).
  \item \amshl{Arousal} $\rightarrow$ systems that initiate and maintain consciousness.
  \item Mediated by reticular activating system (\amshl{RAS}):
    pons $\rightarrow$ midbrain $\rightarrow$ hypothalamus and thalamus.
\end{itemize}
\end{latin}
\begin{itemize}
  \item تغییر حاد هوشیاری \amsarr{} از لتارژی خفیف و کانفیوژن تا کمای عمیق.
  \item بیمار لتارژیک \amsarr{} نمی‌تواند \amshl{\en{arousal state}} خود را حفظ کند.
  \item \amshl{\en{Obtunded}} \amsarr{} \en{$\downarrow$ Arousal, response to stimuli $\downarrow$}.
  \item \amshl{\en{Stupor}} \amsarr{} \en{Only response to pain}.
  \item \en{Coma} \amsarr{} اختلال عملکرد دو نیمکرهٔ مغز و یا در تالاموس و یا ساقهٔ مغز.
\end{itemize}
\subsection*{اختلالات حاد هوشیاری}
شایع‌ترین علل تغییر سطح هوشیاری در کودکان:
\begin{itemize}
  \item مسمومیت.
  \item عفونت‌ها.
  \item تروما به سر \amsarr{} خونریزی مغزی.
  \item هایپوکس ـ ایسکمی \amsarr{} ایست قلبی.
  \item تشنج.
\end{itemize}

\amspage{۲}{pic02.jpg}
\subsection*{ارزیابی کودک با کاهش سطح هوشیاری}
\begin{itemize}
  \item مهم‌ترین علت عوارض طولانی‌مدت در \en{$\downarrow LOC$} \amsarr{} \amshl{هایپوکسی}.
  \item در مواجهه با بیمار \en{LOC} \amsarr{} راه هوایی، تنفس، گردش خون؛ علائم حیاتی، خصوصاً پالس‌اکسی‌متری.
  \item \amshl{برادی‌پنه} \amsarr{} \amshl{\en{$\uparrow ICP$}} / سرکوب‌کننده‌های \en{CNS}.
  \item \amshl{تاکی‌پنه} \amsarr{} هایپوکسی / اسیدوز متابولیک / تب / هرنیشن \amshl{\en{brainstem}}.
  \item فشار خون:
  \begin{itemize}
    \item \amshl{افت} \amsarr{} آسیب تروماتیک مغز (\en{TBI}) / آسیب نخاع.
    \item \amshl{افزایش} \amsarr{} \en{$\uparrow ICP$}.
  \end{itemize}
  \item کودک با \en{LOC} \amsarr{} سطح \amshl{گلوکز} حتماً چک شود.
  \item وضعیت بیمار با تغییر سطح هوشیاری \amsarr{} با \en{GCS} ارزیابی می‌شود.
\end{itemize}

\amspage{۳}{pic03.jpg}
\subsection*{\en{GCS}}
\begin{latin}
\renewcommand{\arraystretch}{1.25}
\begin{center}
\begin{tabular}{|p{12cm}|c|}
\hline
\textbf{Eye opening} & \textbf{Score} \\ \hline
Spontaneous & 4 \\ \hline
To verbal stimuli & 3 \\ \hline
To pain & 2 \\ \hline
None & 1 \\ \hline
\textbf{Verbal} & \\ \hline
Oriented (coos, babbles) & 5 \\ \hline
Confused (irritable, cries) & 4 \\ \hline
Inappropriate words (cries to pain) & 3 \\ \hline
Nonspecific sounds (moans to pain) & 2 \\ \hline
None & 1 \\ \hline
\textbf{Motor} & \\ \hline
Follows commands (spontaneous movement) & 6 \\ \hline
Localizes pain (withdraws to touch) & 5 \\ \hline
Withdraws to pain & 4 \\ \hline
Flexion to pain (abnormal flexion) & 3 \\ \hline
Extension to pain (abnormal extension) & 2 \\ \hline
None & 1 \\ \hline
\end{tabular}
\end{center}
\end{latin}
\begin{itemize}
  \item اختلال در پاسخ‌های مردمک:
  \begin{itemize}
    \item \en{fix} شدن \amsarr{} آسیب \en{midbrain}.
    \item \amshl{\en{pinpoint}} \amsarr{} آسیب \amshl{\en{pontine}}.
  \end{itemize}
  \item در بیماران کمایی، حرکات چشم با \en{Oculocephalic vestibular reflexes}:
  \begin{itemize}
    \item \amshl{\en{Doll's-eye maneuver}}.
    \item اگر غیرایمن بود \amsarr{} \en{Oculovestibular response} \amsarr{} آب سرد داخل گوش.
  \end{itemize}
\end{itemize}

\amspage{۴}{pic04.jpg}
\begin{itemize}
  \item اگر ساقهٔ مغز (\en{brainstem}) سالم باشد \amsarr{} چشم به سمت گوش تحریک‌شده می‌چرخد.
\end{itemize}
\subsection*{\en{Posturing}}
\begin{latin}
\begin{itemize}
  \item Decorticate:
  \begin{itemize}
    \item Rigid extension of lower extremities.
    \item Flexion and supination of the arms.
    \item Fisting of the hands.
  \end{itemize}
  \item Decerebrate:
    \amshl{Rigid extension and internal rotation} of the upper and lower extremities.
\end{itemize}
\end{latin}
\begin{itemize}
  \item دکورتیکه \amsarr{} ساقهٔ مغز سالم است.
  \item دسربره \amsarr{} آسیب مغز میانی / \amshl{اختلالات توکسیک متابولیک}.
  \item این پوسچرها می‌توانند قرینه یا غیرقرینه باشند:
    توکسیک متابولیک / اینتراکرانیال.
\end{itemize}
\subsection*{اتیولوژی}
\begin{itemize}
  \item دسته‌بندی کلی علت‌های کما:
  \begin{itemize}
    \item توکسیک ـ متابولیک.
    \item عفونت.
    \item ساختمانی \amsarr{} افت هوشیاری ناگهانی.
  \end{itemize}
  \item عفونت \en{CNS} \amsarr{} افت هوشیاری ناگهانی / مننگوانسفالیت \en{HSV} \amsarr{} تحت‌حاد.
  \item علائم تحریک مننژ \amsarr{} فوتوفوبی، برودزنسکی، درد حین حرکت و خم کردن گردن.
  \item سردرد + تهوع جهنده‌ای \amsarr{} علائم \en{$\uparrow ICP$}.
\end{itemize}

\amspage{۵}{pic05.jpg}
\subsection*{افزایش \en{ICP} و بررسی بیمار \en{AMS}}
\begin{itemize}
  \item افزایش \en{ICP} \amsarr{} ادم پاپیل، فلج عصب \en{III}، انحراف چشم به پایین و خارج؛ فلج عصب \en{VI} (ناتوانی حرکت چشم به خارج).
  \item تریاد کوشینگ \amsarr{} \amshl{برادی‌کاردی + برادی‌پنه + هایپرتنشن} (\amshl{\en{$\uparrow ICP$}}).
  \begin{itemize}
    \item هرنی از فورامن مگنوم \amsarr{} توقف قلب و تنفس.
  \end{itemize}
  \item تصویربرداری \amsarr{} در تمامی موارد اورژانس بررسی بیمار \en{AMS}:
    \en{CT} بدون کنتراست.
  \item حضور \en{RBC} در \en{LP} \amsarr{} خونریزی \en{SAH}، خونریزی پارانشیم، عفونت ویروسی (\en{HSV}) دهندهٔ خونریزی.
  \item کودک با سنکوپ بدون علت \amsarr{} بررسی کامل قلبی.
  \item \en{$\frac{2}{3}$} کودکان با سنکوپ دچار \amshl{\en{Anoxic seizure}} می‌شوند \amsarr{} \en{DDX} صرع.
\end{itemize}
\subsection*{سودوتومور سربری — \en{Idiopathic intracranial hypertension}}
\begin{itemize}
  \item تصویربرداری \en{NL} ولی \en{$\uparrow ICP$}.
  \item سردرد، وزوز، فلج عصب \en{VI}، ادم پاپیل.
  \item ناشی از برخی داروها \amsarr{} \amshl{تتراسایکلین، \en{Vit A}، \en{OCP}}، ؟
  \item بیماری آدیسون، بیماری‌های تیروئید.
  \item افزایش وزن ناگهانی ریسک فاکتور است.
  \item درمان: \amshl{استازولامید، دیورتیک، توپیرامات}.
\end{itemize}
\begin{itemize}
  \item سردرد در \en{$\uparrow ICP$} \amsarr{} با دراز کشیدن (خواب) بدتر می‌شود؛ با سرفه بدتر می‌شود: \amshl{مانور والسالوا}.
  \item اندیکاسیون انتوباسیون در \en{$\uparrow ICP$}:
    هایپوکسی راجعه، هایپوونتیلیشن، \en{$GCS\leq 8$}، از دست رفتن رفلکس‌های محافظ راه هوایی.
\end{itemize}
\begin{latin}
Slight hyperventilation to a $PaCO_2 \geq 35\,\mathrm{mmHg}$ $\longrightarrow$ $\downarrow ICP$.
\end{latin}

\amspage{۶}{pic06.jpg}
\subsection*{کانکاشن و پروتکل بازگشت}
\begin{itemize}
  \item پروتکل بازگشت به ورزش پس از کانکاشن: هر گام \en{24 h}.
\end{itemize}
\begin{enumerate}
  \item بازگشت به مدرسه.
  \item فعالیت هوازی خفیف.
  \item ورزش.
  \item تمرین.
  \item تمرین.
  \item مسابقه.
\end{enumerate}
% Steps 3--5 retain the source's abbreviated labels (ورزش، تمرین، تمرین).
\begin{itemize}
  \item کانکاشن غیرعارضه‌دار نیازی به \amshl{تصویربرداری ندارد}؛ اگر شک داریم ضایعهٔ ساختمانی باشد.
  \item سردرد پس از ضربه خیلی \amshl{شایع است}؛ اگر هوشیاری \en{full} شود نیاز به تصویربرداری نیست.
  \item اگر بعد از ۱ تا ۲ روز علائم باقی ماند \amsarr{} تصویربرداری.
  \item شکستگی دپرشن \amsarr{} اگر بالای \en{0.5--1 cm} \amsarr{} جراحی.
  \item شکستگی‌های خطی، دیاستاتیک، دپرشن \amsarr{} \amshl{حتماً تصویربرداری}.
  \item هرنی ساب‌فالکس \amsarr{} ضعف، \amshl{بی‌اختیاری ادرار}، کما.
\end{itemize}

\ifnum\AMSStandalone=1 \expandafter\enddocument\fi
